\documentclass[12pt]{amsart}
\usepackage[all]{xy} % for commutative diagrams
\usepackage{amsmath} % various goodies
\usepackage{amsthm} % some stuff for theorem numbering
\usepackage{amsfonts}

%% Theorem-type environments.
% The first parameter to \newtheorem is what the environment is called in the 
% LaTeX source. (What goes into the braces after \begin and \end.)
% The second is the name that shows up in the text.

\theoremstyle{plain} 
\newtheorem{theorem}{Theorem}
\newtheorem{lemma}{Lemma}
\newtheorem{corollary}{Corollary}
\newtheorem{proposition}{Proposition}
\newtheorem{conjecture}{Conjecture}
% (Criterion, Algorithm)

\theoremstyle{definition} 
\newtheorem{definition}{Definition}
\newtheorem{example}{Example}
% (Condition, Problem)

\theoremstyle{remark} 
\newtheorem{remark}{Remark}
\newtheorem{note}{Note}
% (Notation, Claim, Summary, Acknowledgement, Case, Conclusion)

%% Math operators (easy thanks to amsmath)
\DeclareMathOperator{\Aut}{Aut}
\DeclareMathOperator{\core}{Core}
\DeclareMathOperator{\orb}{Orb}
\DeclareMathOperator{\cay}{Cay}
\DeclareMathOperator{\sab}{Sab}

% Commands:
\newcommand{\cB}{\mathcal{B}}
\newcommand{\cS}{\mathcal{S}}
\newcommand{\cP}{\mathcal{P}}
\newcommand{\cZ}{\mathcal{Z}}
\newcommand{\cD}{\mathcal{D}}
\newcommand{\cH}{\mathcal{H}}
\newcommand{\cQ}{\mathcal{Q}}
\newcommand{\cK}{\mathcal{K}}
\newcommand{\cN}{\mathcal{N}}

% Lazy:
\newcommand{\travH}{\{ \tau_1, \tau_2, \dots \tau_n \}}
\newcommand{\travR}{\{ t_1, t_2, \dots t_b \}}
\newcommand{\rangen}{\{1 \dots n \}}

\begin{document}

\section{$G$-symmetric designs}

In this Chapter we shall define a new category $G$-Design.
The objects of this category are $G$-symmetric designs
and the morphisms are $G$-design homomorphisms.
We shall exhibit a ``forgetful functor'' from $G$-graph
to $G$-design and rephrase some of the questions we have
been asking about $G$-symmetric graphs into questions 
about $G$-symmetric designs.
In particular we shall describe a ``decomposition''
for $G$-symmetric designs that is analogous to the 
decomposition given by Gardiner and Praeger for $G$-symmetric graphs.

\begin{definition}
A $G$-flag transitive design $\cD$ is {\it self-dual}
if there exists a $G$-isomorphism $\rho = (\rho_P, \rho_B)$
between $\cD$ and $\cD^{*}$. 
The $G$-isomorphism $\rho$ is called a {\it duality} of $\cD$.
\end{definition}

\begin{definition}
A {\it $G$-symmetric design} is self-dual $G$-flag transitive design $\cD$
which admits a dulaity $\rho = (\rho_P, \rho_B)$ with the property that 
$\rho_B \circ \rho_P = id_P$ and $\rho_P \circ \rho_B = id_B$.
In this case the $G$-isomorphism $\rho$ is called a {\it polarity}
of $\cD$.
\end{definition}

Observe that if $\cD$ is a $G$-symmetric design with the property
that $(p, \rho(p))$ is a a flag for some $p \in P_\cD$,
then since $\rho$ is a $G$-isomorphism it follows
that the point-stabilizer of $\cD$ is isomorphic to the flag-stabilizer of $\cD$.
Thus each point is incident with exactly one block
and vice versa. 
We consider such $G$-symmetric designs to be ``degenerate'',
and unless an explicit statement to the contrary is given shall
take the expression ``$G$-symmetric design'' 
to mean ``non-degenerate $G$-symmetric design''.

In a sense that will become clear a little later, 
these ``degenerate'' $G$-symmetric
designs correspond to the ``degenerate'' orbital graph $\orb_\Delta(G,\Omega)$
which could be formed from the permutation group $(G, \Omega)$
by taking $\Delta$ to be the diagonal orbit,
and also to the ``degenerate'' Sabidussi graph $\sab(G,H,a)$
which could be formed by taking $a \in H$.
\hfill \break

We shall see that each $G$-symmetric graph 
gives rise in a natural way to a $G$-symmetric design 
and conversely a $G$-symmetric design together 
with a ``marked'' polarity give rise to a $G$-symmetric graph.
We shall also see that in some cases, by choosing a different
polarity it is possible to construct two non-isomorphic $G$-symmetric
graphs from the same $G$-symmetric design.
For any graph $\Gamma$, let $N \Gamma = \{ \Gamma(v) : v \in V \Gamma \}$
denote the set of {\it neighbourhoods} of $\Gamma$.

\begin{proposition}
If \/ $\Gamma$ is a $G$-symmetric graph,
then the incidence structure $\cD(\Gamma) = (V \Gamma, N \Gamma,I)$
where $(v,n) \in I$ if and only if $v \in n$ is a $G$-symmetric design
\begin{proof}
The $G$-arc transitivity of $\Gamma$ is sufficient to ensure
that $\cD(\Gamma)$ is a $G$-flag transitive design.
To see that it is in fact
a $G$-symmetric design we must exhibit a polarity.

Let $\rho_P : P \to B$ be given by $v \mapsto \Gamma(v)$.
The map $\rho_P$ is clearly bijective, and since $G$ acts on $\Gamma$
as a group of automorphisms we have:
\begin{eqnarray*}
\rho_P(v^g) & = & \Gamma(v^g) \\
            & = & \Gamma(v)^g \\
	    & = & \rho_P(v)^g.
\end{eqnarray*}
So $\rho_P$ induces a permutation isomorphism.
Now, take $\rho_B = \rho_P^{-1}$.
We must check that the pair $\rho = (\rho_P, \rho_B)$ preserve
the incidence structure of the design.

Suppose that $(v,\Gamma(w))$ is a flag of $\cD$, so $v \in \Gamma(w)$. 
Since $\Gamma$ is a simple graph it follows immediately that $w \in \Gamma(v)$
and so $(w, \Gamma(v))$ is also a flag of $\cD$.
That is $(\Gamma(v), w)$ is a flag of $\cD^{*}$.
But 
\begin{eqnarray*}
\rho((v, \Gamma(w))) & = & (\rho_P(v), \rho_B(\Gamma(w))) \\
                     & = & (\Gamma(v), w).
\end{eqnarray*}
So we are done.

\end{proof}
\end{proposition}

\begin{proposition}
If $\cD$ is a (non-degenerate) $G$-symmetric design with ``marked'' polarity $\rho$
then the graph $\Gamma(\cD, \rho)$ with vertex set $P_\cD$ 
and arc set $\{ (p,q) : (q,\rho_P(p)) \in I_\cD \}$ is $G$-symmetric. 
\begin{proof}
We must first check that $\Gamma(\cD, \rho)$ is well-defined.
The non-degenerateness of $\cD$ ensures that there are no loops. 
If $(p,q)$ is an arc of $\Gamma(\cD, \rho)$
then $(q, \rho_P(p))$ is a flag of $\cD$.
Since $\rho$ is an isomorphism, 
if $(q, \rho_P(p))$ is a flag of $\cD$
then $(\rho_P(q), \rho_B \circ \rho_P(p)) = (\rho(q), p)$
is a flag of $\cD^{*}$.
It follows that $(p,\rho(q))$ is a flag of $\cD$,
and thus $(q,p)$ is an arc of $\Gamma(\cD, \rho)$.
So $\Gamma(\cD, \rho)$ is simple.

By proposition [?] $G$ acts transitively on the points $\cD$
so $\Gamma(\cD, \rho)$ is $G$-vertex transitive.
Suppose that $(p,q_1)$ and $(p,q_2)$ are two distinct arcs of
$\Gamma(\cD, \rho)$. Then we have $q_1, q_2 \in T(\rho(p))$.
By proposition [?] 
the stabilizer of $\rho(p)$ acts transitively on $T(\rho(p))$.
Thus we can find some $g \in G$ which fixes $\rho(p)$ and 
carries $q_1$ to $q_2$. 
Since $\rho$ is a $G$-isomorphism, if $g$ fixes $\rho(p)$
it also fixes $p$, thus $g$ carries the arc $(p,q_1)$
to the arc $(p,q_2)$. So $\Gamma(\cD, \rho)$ is $G$-locally
transitive. The result follows.

\end{proof}
\end{proposition}

\begin{proposition}
For any design $\cD$ and any polarity $\rho$ the design $\cD(\Gamma(\cD,\rho))$
is isomorphic to $\cD$
\begin{proof}
later
\end{proof}
\end{proposition}

\begin{definition}
If $\cD_1$ and $\cD_2$ are two $G$-symmetric designs
then a $G$-design homomorphism from $\cD_1$ to $\cD_2$
is a pair of permutation homomorphisms $\eta_P : P_1 \to P_2$ and $\eta_B : B_1 \to B_2$
satisfying $(p,b) \in I_1$ implies $(\eta_P(p),\eta_B(b)) \in I_2$. 
\end{definition}

\begin{proposition}
If $\cD_1$ and $\cD_2$ are two $G$-symmetric designs
then a $G$-design homomorphism from $\cD_1$ to $\cD_2$
induces a a $G$-permutation homomorphism from the flags of $\cD_1$ 
to the flags of $\cD_2$
\end{proposition}

What I want to say is that $G$-symmetric designs
together with $G$-design homomorphisms form a category
(well, kind of more of a poset really).
I also want to say that the category $G$-Graph
``projects down'' onto the category $G$-Design.
$F : G$-Graph $\to G$-Design.
If there's a $G$-graph homomorphism from $x$ to $y$
then there's a $G$-design homomorphism from $F(x)$ to $F(y)$

\section{Kernel of a design homomorphism}

Gardiner and Praeger showed that whenever we have
a $G$-symmetric graph homomorphism $\Gamma \mapsto \Sigma$
we get an induced ``cross-sectional'' design
on each of the fibers of the kernel.
Here I'll show that you get exactly the same thing
whenever you have $G$-symmetric design homomorphism.

So, with an abuse of notation, let $\pi : \cD \to \cQ$
be a $G$-symmetric design homomorphism.
For any point $q \in P_\cQ$
consider the {\it fiber} over $q$, 
that is the set $F(q) = \pi^{-1}(q) = \{ p \in P_\cD : \pi(p) = q \}$.
The ``induced kernel design'' on the fiber of $q$
is $\cK(q) = (F(q), T(q), I_q)$
where $(p,d) \in I_q$ if and only if there is some $b \in \pi^{-1}(d)$
such that $(p,b) \in I_\cD$. 

I think that if $\Gamma \mapsto \Sigma$ is a $G$-symmetric graph
homomorphism and $\cD$ is the ``GP-cross-section'' of the map
then if $F(\Gamma) \mapsto F(\Sigma)$ is the $G$-symmetric design
homomorphism then the 

%Define an incidence relation $I_p \subseteq F(\bar{p}) \times B(\bar{p})$
%as follows: $(p,b) \in I_{\bar{p}}$ if and only if
%there exists some $b \in \pi_B^{-1}(\bar{b})$ 
%such that $(p,b) \in I$.
%
%\begin{proposition}
%If $H$ is the stabilizer of $\bar{p}$,
%then $D_{\bar{p}} = (F(\bar{p}), B(\bar{p}), I_{\bar{p}})$ is a $H$-design.
%
%\begin{proof}
%Suppose that $q \in F(\bar{p})$ and $d \in B(\bar{p})$
%such that $(q,d) \in I_{\bar{p}}$.
%Then there exists some $b \in \pi_B(d)$ such that
%$(q,b) \in I$. Since $H$ acts on $D$ as a group of
%automorphisms, it follows that for any $h \in H$
%we also have $(q^h, b^h) \in I$.
%Since $b^h \in \pi_B(d)^h = \pi_b(d^h)$
%we have $(q^h, d^h) \in I_{\bar{p}}$,
%so $H$ acts on $D_{\bar{p}}$
%as a group of automorphisms.
%
%The set $F(\bar{p})$ is a subset of $P$,
%in fact it is a block of imprimitivity of $P$ with setwise stabilizer $H$.
%Since $G$ acts transitively on $P$ for any
%$q_1, q_2 \in F(\bar{p})$ we can find some $g \in G$ such that
%$q_1^g = q_2$. Since $F(\bar{p})$ is a block of imprimitivity
%we must have  $F(\bar{p})^g = F(\bar{p})$, so in fact $g \in H$.
%That is $H$ acts transitively on $F(\bar{p})$.
%
%Let $\bar{b} = \rho_P(\bar{p})$ and let $P(\bar{b})$
%denote the set of points incident with $\bar{b}$.
%
%Let $d_1, d_2 \in B(\bar{p})$ be any two blocks of $D_Q$.
%Since $G$ acts transitively on the blocks of $D_Q$,
%we can find some $g \in G$ such that $d_1^g = d_2$.
%
%\end{proof}
%\end{proposition}
%
%We call the design design $(F(\bar{p}), B(\bar{p}), I_{\bar{p}})$
%the kernel design induced on the fiber of $\bar{p}$.
%
%\begin{proposition}
%Upto isomorphism $(F(\bar{p}), B(\bar{p}), I_{\bar{p}})$
%is independent of the choice of $\bar{p}$.
%\begin{proof}
%\end{proof}
%\end{proposition}
%
%
%\begin{theorem}
%This kernel design thing is, in fact, the design
%Gardiner and Praeger were talking about in there paper.
%\begin{proof}
%maybe?
%\end{proof}
%\end{theorem}
%
%
%Let $p \in P$ be any point of $\cD$ and let $b = \rho(p)$.
%If $H < G$ stabilizes $p$ then $H$ also stabilizes $b$,
%thus the permutation groups $(G,P)$ and $(G,B)$ are both
%permutation isomorphic to $(G, \cos(H))$ and thus also to each other.

%Suppose that $d \in B$ is a block incident with $p$, 
%that is, $(p,d)$ is a flag of $\cD$.
%Since $\cD$ is a $G$-design we can find some $g \in G$ such that
%$b^g = d$. It follows that the stabilizer of $d$ is $g^{-1}Hg$,
%and thus the stabilizer of the flag $(p,d)$ is $R = H \cap g^{-1}Hg$.
%The action of $H$ on $B(p)$ is permutation isomorphic to $(H, \cos(R))$
%and the action of $g^{-1}Hg$ on $P(b)$ is permutation isomorphic to $(g^{-1}Hg, \cos(R))$.
%These two actions are permutationally equivalent but not isomorphic
%unless $d = b$ in which case each point of the design is incident with a unique block
%and the design is somewhat trivial.


%Suppose that $D = (P,B,I)$ is a $G$-symmetric design,
%so $(G,P)$ is a transitive permutation group.
%Suppose further that $\cP$ is a $G$-invariant partition of P.
%Let $\pi_P : P \to \cP$ be the natural projection
%which sends a point to the block of the partition containing it.
%The partition $\cP$ gives rise to an equivalence relation on $P$,
%namely $p_1 \equiv p_2$ iff $\pi_P(p_1) = \pi_P(p_2)$.
%If $\rho$ is some polarity of $D$ then we have an induced
%equivalence relation on the blocks of $D$ given by
%$b_1 \equiv b_2$ iff $\pi_P \circ \rho^{-1}(b_1) = \pi_P \circ \rho^{-1}(b_2)$.
%Let $\cB$ denote the equivalence classes of this relation
%and let $\pi_B$ be the natural projection from $B$ into $\cB$.
%Define a new design $\overline{D} = (\cP,\cB,\overline{I})$
%where the incidence relation is
%$(\bar{p},\bar{b}) \in \overline{I}$ if and only if there is
%some $p \in P$ and some $b \in B$ such that $\bar{p} = \pi_P(p)$
%and $\bar{b} = \pi_B{b}$ and $(p,b) \in I$.
%
%\begin{theorem}
%The design $\overline{D} = (\cP,\cB,\overline{I})$
%is $G$-symmetric
%The maps $\pi_P$ and $\pi_B$ induce a $G$-epimorphism
%from $D$ to $\overline{D}$. Furthermore, every $G$-epimorphism
%is of this form
%\begin{proof}
%obvious?
%\end{proof}
%\end{theorem}
%
\end{document}
